\chapter{Library Architecture}

\section{Module Overview}

\begin{table}[h]
\centering
\begin{tabular}{lll}
\hline\hline
\textbf{Module} & \textbf{Purpose} & \textbf{Key Functions} \\
\hline
\texttt{domain.h/c} & Data structures & \texttt{LatticeInfo}, \texttt{Star}, \texttt{SAFBBasis} \\
\texttt{symmetry\_ops.h/c} & Symmetry operations & \texttt{read\_spacegroup\_ops\_txt}, \texttt{star\_from\_hkl} \\
\texttt{basis.h/c} & Basis construction & \texttt{family\_planes\_info}, \texttt{build\_basis} \\
\texttt{initializers.h/c} & Amplitude assignment & \texttt{random\_initialization}, \texttt{manual\_initialization} \\
\texttt{field.h/c} & Field generation & \texttt{build\_field}, \texttt{build\_field\_2d} \\
\texttt{engine.h/c} & High-level API & \texttt{engine\_create}, \texttt{engine\_init\_random} \\
\texttt{analytic.h/c} & Analytical calculations & \texttt{calculate\_square\_norm} \\
\hline\hline
\end{tabular}
\caption{Library modules and their responsibilities}
\end{table}

\section{Data Structures}

\subsection{LatticeInfo}

\begin{lstlisting}
typedef struct {
    double Lx, Ly, Lz;      // Lattice lengths
    double alpha, beta, gamma; // Lattice angles (degrees)
    double G[3][3];         // Reciprocal metric tensor
    double Ginv[3][3];      // Inverse metric tensor
    int dim;                // 2 or 3
} LatticeInfo;
\end{lstlisting}

\subsection{Star}

\begin{lstlisting}
typedef struct {
    int hkl[3];             // Canonical Miller indices
    double q2;              // |q|^2 (q-squared)
    int multiplicity;       // Number of vectors
    int closed;             // 1 if closed under inversion
    int vectors[3][MAX_STAR_SIZE]; // Symmetry-equivalent vectors
    int vector_keys[MAX_STAR_SIZE]; // Lexicographic keys
    double phase_factors[MAX_STAR_SIZE]; // Phase factors
} Star;
\end{lstlisting}

\subsection{SAFBBasis}

\begin{lstlisting}
typedef struct {
    int n_modes;            // Number of modes (families)
    Star modes[MAX_MODES];  // Mode data
    int centrosymmetric;    // 1 if inversion is in point group
    int inversion_at_origin; // 1 if inversion at origin
} SAFBBasis;
\end{lstlisting}

\section{API Usage}

\begin{lstlisting}
// Create engine from space group file
Engine *engine = engine_create("examples/space_groups_3d/ia_3d.txt",
                                "examples/lattice_cubic.txt");

// Random initialization
engine_init_random(engine, 42);

// Build field
engine_build_field(engine, 64, 64, 64);

// Output field to VTK
engine_output_field(engine, "gyroid.vts");

// Clean up
engine_destroy(engine);
\end{lstlisting}
